\documentclass[10pt,a4paper]{amsart}
\usepackage[margin=19mm]{geometry}
\usepackage{amsmath,amssymb,mathtools}
\usepackage[scr=rsfs,cal=euler]{mathalfa}
\usepackage{xfrac}
\usepackage[unicode,hidelinks,bookmarks=false]{hyperref}
\newcommand{\ie}{i.e.\ }
\newcommand{\cf}{cf.\ }
\newcommand{\resp}{respectively\ }
\newcommand{\Subsection}{\subsection}
\providecommand{\nobreakdash}{\nobreak\hbox{-}}
\setlength{\emergencystretch}{2em}
\multlinegap0pt
\usepackage{orcidlink}
\usepackage{stmaryrd,url,float}
\renewcommand{\labelenumi}{\theenumi}
\newcommand*\fixitem {\item[]%
  \refstepcounter{enumi}\hskip-\leftmargin\labelenumi\hskip\labelsep}
 \newtheorem{theorem}{Theorem}[section]
 \newtheorem{lemma}[theorem]{Lemma}
  \newtheorem{proposition}[theorem]{Proposition}
 \newtheorem{corollary}[theorem]{Corollary}
 \newtheorem{remark}[theorem]{Remark}
 \newtheorem*{special1}{Named Theorem}
 \newtheorem*{special2}{Theorem}
\def\Star{\ensuremath{{\bf(*)}}}
\def\DCP{\ensuremath{{\bf(DCP_{A})}}}
\def\CEP{\ensuremath{{\bf(CEP)}}}
\def\iDEP{\ensuremath{{\bf(DEP1)}}}
\def\iiDEP{\ensuremath{{\bf(DEP2)}}}
\def\HDP{\ensuremath{{\bf(HDP)}}}
\def\JDP{\ensuremath{{\bf(JDP)}}}
\def\LDP{\ensuremath{{\bf(LDP)}}}
\def\RNRP{\ensuremath{{\bf(RNRP)}}}
\def\RRP{\ensuremath{{\bf(RRP)}}}
\def\iLTP{\ensuremath{{\bf(LP1)}}}
\def\R{\mathbb{R}}
\def\F{\mathbb{F}}
\def\P{\mathbb{P}}
\def\N{\mathbb{N}}
\def\AP{\ensuremath{{\bf(AP)}}}
\def\CC{\ensuremath{{\bf(CC)}}}%Countable Cofinality
\def\sNIP{\ensuremath{{\bf(NIP_{\text{s}})}}}
\def\wNIP{\ensuremath{{\bf(NIP_{\text{w}})}}}
\def\MCP{\ensuremath{{\bf(MCP)}}}
\def\FCP{\ensuremath{{\bf(FCP)}}}
\def\DCP{\ensuremath{{\bf(DCP)}}}
\def\DFP{\ensuremath{{\bf(DFP)}}}
\def\CPP{\ensuremath{{\bf(CPP)}}}
\def\LIP{\ensuremath{{\bf(LIP)}}}
\def\LMP{\ensuremath{{\bf(LMP)}}}
\def\ES{\ensuremath{{\bf(ES)}}}
\def\CA{\ensuremath{{\bf(CA)}}}
\def\lbrak{\left[}
\def\rbrak{\right]}
\def\lpar{\left(}
\def\rpar{\right)}
\def\abClosed{\lbrak a,b\rbrak}
\def\stepLatt{\mathcal{L}}
\def\elemLatt{\stepLatt(\abClosed)}
\def\elemI{I_{\abClosed}}
\def\istepLatt{\stepLatt_u}
\def\into{\longrightarrow}
\def\sub{\subseteq}
\def\setdiff{\backslash}
\def\nseqlim{{\displaystyle\lim_{n\rightarrow\infty}}}
\def\nseqliminf{{\displaystyle\liminf_{n\rightarrow\infty}}}
\def\kseqlim{{\displaystyle\lim_{k\rightarrow\infty}}}
\def\abRing{\mathfrak{S}}
\def\abAlg{\mathcal{A}}
\def\genSub{\mathcal{S}}
\def\genSubAlg{\left<\genSub\right>}
\def\genInt{\mathcal{I}}
\def\genIntAlg{\left<\genInt\right>}
\def\Cara{Carath\'{e}odory\  }
\usepackage{tikz}
\usetikzlibrary{arrows.meta, positioning}
\begin{document}
\title[Extension Theorems in Measure and Integration Theory as Comp]{Extension Theorems in Measure and Integration Theory as Completeness Axioms}
\author{Rafael Reno S. Cantuba}
\date{}
\maketitle
\noindent\emph{Source and licence.} Extension Theorems in Measure and Integration Theory as Completeness Axioms. Version arXiv:2609.29510v1 (CC BY 4.0). This material is distributed under the Creative Commons Attribution 4.0 International licence, http://creativecommons.org/licenses/by/4.0/, as recorded in the arXiv OAI licence metadata for this exact version and shown on the abstract page https://arxiv.org/abs/2609.29510v1. Full rights notice: licenses/arxiv-2609.29510-CC-BY-4.0.txt.\par\smallskip
\noindent\textit{Editorial scope.} Counted unit: the original \texttt{theorem} environment \texttt{-} at source lines 328--329 together with its complete proof at lines 330--338; the delivered document keeps source lines 214--339 so that every referenced label and every citation resolves inside it. Native editable source is the paper's own concatenated TeX; the AMS wrapper is editorial formatting only. No publisher class, upstream script or unrelated artwork is redistributed.
\section{Statements of the Dedekind completeness axioms}\label{StateSec}

We are now ready to give the statements of the measure-theoretic and integration-theoretic axioms for $\F$ that we shall prove to be equivalent.

\begin{itemize}
\item[\CPP] \emph{Carath\'{e}odory Premeasure Property.} Given $\abClosed\sub\F$, the function\linebreak $E\mapsto\int_a^b\chi_E$ is a premeasure on the algebra of all measurable subsets of $\abClosed$.
    \item[\CEP]\emph{Carath\'{e}odory Extension Property.} Given $\abClosed\sub\F$, the function\linebreak $E\mapsto\int_a^b\chi_E$ is extendable to a \Cara outer measure $\mu^*$, and the collection of all $\mu^*$-measurable subsets of $\abClosed$ is a $\sigma$-algebra that contains the algebra of all measurable subsets of $\abClosed$.
    \item[\DFP]\emph{Daniell Functional Property.} Given $\abClosed\sub\F$,  the positive linear functional $\varphi\mapsto\int_a^b\varphi$ is a Daniell functional on the vector lattice of all step functions $\abClosed\into\F$.
    \item[\iDEP]\emph{First Daniell Extension Property.} Given $\abClosed\sub\F$,  the positive linear functional $I_{\abClosed}:\varphi\mapsto\int_a^b\varphi$ on the vector lattice $\stepLatt(\abClosed)$ of all step functions $\abClosed\into\F$ has a continuous extension to the first Daniell extension $\istepLatt(I_{\abClosed})$ of $\stepLatt(\abClosed)$. 
\item[\iiDEP]\emph{Full Daniell Extension Property.} Given $\abClosed\sub\F$, the positive linear functional $I_{\abClosed}:\varphi\mapsto\int_a^b\varphi$ on the vector lattice $\stepLatt(\abClosed)$ of all step functions $\abClosed\into\F$ is extendable to a Daniell integral on the second Daniell extension $L^1(\stepLatt(\abClosed),I_{\abClosed})$ of $\stepLatt(\abClosed)$. 

    \item[\MCP]\emph{Monotone Convergence Property.} Given $\abClosed\sub\F$, for each monotonically increasing sequence $(\varphi_n)$ of nonnegative Lebesgue measurable functions $\varphi_n:\abClosed\into\F$ that converges pointwise to $f:\abClosed\into\F$ almost everywhere on $\abClosed$, we have $\int_a^bf=\nseqlim\int_a^b\varphi_n$.
    
    \item[\FCP]\emph{Fatou Convergence Property.} Given $\abClosed\sub\F$, for each sequence $(\varphi_n)$ of nonnegative Lebesgue measurable functions $\varphi_n:\abClosed\into\F$ that converges pointwise to $f:\abClosed\into\F$ almost everywhere on $\abClosed$, the quantity $\nseqliminf\int_a^b\varphi_n$ exists in $\F$. Furthermore, $\int_a^bf\leq\nseqliminf\int_a^b\varphi_n$.
    \item[\DCP]\emph{Dominated Convergence Property.} Given $\abClosed\sub\F$, for each sequence $(\varphi_n)$ of Lebesgue measurable functions $\varphi_n:\abClosed\into\F$ that converges pointwise to $f:\abClosed\into\F$ almost everywhere on $\abClosed$, if there exists a Lebesgue measurable function $g:\abClosed\into\F$ such that $|\varphi_n|\leq g$ almost everywhere on $\abClosed$ for all $n\in\N$, then $\int_a^bf=\nseqlim\int_a^b\varphi_n$.


    \item[\LIP]\emph{Lebesgue Integral Property.} Given $\abClosed\sub\F$, every Lebesgue measurable function has a Lebesgue integral.

        \item[\LMP]\emph{Lebesgue Measure Property.} Given $\abClosed\sub\F$, every  measurable subset of $\abClosed$ has Lebesgue measure.

   % \item[\iLTP]\emph{Littlewood's First Principle.} Given $\abClosed\sub\F$, for each measurable set $E\sub\abClosed$ and each $\varepsilon\in\P$, there exist intervals $I_1,I_2,\ldots,I_n\sub\abClosed$ such that if $U=\bigcup_{i=1}^nI_i$, then $\int_a^b\chi_{E\Delta U}<\varepsilon$.
\end{itemize}

Each of the above statements has a general form \Star, which is the statement that results when we replace $\abClosed$ by an arbitrary subset $X$ of $\F$, or the function $E\mapsto\int_a^b\chi_E$ by an arbitrary measure $\mu$ on $X$, or $\elemI$ by an arbitrary positive linear functional $I$ on an arbitrary vector lattice on $X$, and so on. This general form \Star\  is a theorem for the case $\F=\R$, which is the unique Dedekind complete field up to isomorphism, so \Star\  follows from Dedekind completeness, which implies the specific case, which is, say, when $X=\abClosed$, $\mu:E\mapsto\int_a^b\chi_E$, $I=\elemLatt$, and so on. Thus, our goal, similar to that done in \cite[pp.~116--117]{can24}, is to show that the specific forms of the statements, the ones written above, imply Dedekind completeness of $\F$, and this completes the circle of equivalence. In particular, \LIP\  and \LMP, together with Littlewood's Three Principles, have already been proven to be equivalent to Dedekind completeness in \cite{can24}, so we have the more specific goal of proving that each of the other statements is a sufficient condition for \LIP\  or \LMP. We refer the reader to \cite{can24} for a full discussion of the involvement of Littlewood's Three Principles in reverse mathematics. 

\section{Logical relationships between the statements}

The general plan for proving the equivalence of the statements in Section~\ref{StateSec} with the Dedekind completeness of $\F$ is summarized in the following figure.

\begin{figure}[H]
\centering
\begin{tikzpicture}[
scale=0.9,             % Scales the coordinates (70% size)
    transform shape,       % Scales the nodes/shapes to match
    double arrow/.style={
        double, 
        double distance=1.5pt, 
        draw=#1, 
        ->, 
        >={Latex[length=4pt, width=5pt, color=#1]}, 
        shorten >= 3pt, 
        shorten <= 3pt,
        rounded corners=10pt % This rounds the "corner" of the edge itself
    },
    double iffarrow/.style={
        double, 
        double distance=1.5pt, 
        draw=#1, 
        <->, 
        >={Latex[length=4pt, width=5pt, color=#1]}, 
        shorten >= 3pt, 
        shorten <= 3pt,
        rounded corners=10pt % This rounds the "corner" of the edge itself
    },
    vertex/.style={
        draw, 
        rectangle, 
        rounded corners=5pt, 
        minimum width=40pt, 
        minimum height=20pt
    }
]

        \node[vertex, draw=black, fill=white] (iDEP) at (0,15) {\iDEP};
        \node[vertex, draw=black, fill=white] (DFP) at (0,25.5*0.5) {\DFP};
        \node[vertex, draw=black, fill=lightgray] (LIP) at (0,10.5) {\LIP};

        \node[vertex, draw=black, fill=lightgray] (iiDEP) at (3.875,15) {\iiDEP};
        \node[vertex, draw=black, fill=lightgray] (CEP) at (3.875,13.5) {\CEP};
        \node[vertex, draw=black, fill=white] (CPP) at (3.875,12) {\CPP};
        \node[vertex, draw=black, fill=lightgray] (LMP) at (3.875,10.5) {\LMP};

        \node[vertex, draw=black, fill=white] (MCP) at (2*3.875,15) {\MCP};
        \node[vertex, draw=black, fill=white] (FCP) at (2*3.875,25.5*0.5) {\FCP};
        \node[vertex, draw=black, fill=white] (DCP) at (2*3.875,10.5) {\DCP};

     \draw[double arrow=black] (iiDEP) -- (iDEP);
    \draw[double arrow=black] (iDEP) -- (DFP);
    \draw[double arrow=black] (DFP) -- (LIP);
    \draw[double iffarrow=black] (LIP) -- (LMP);

    
    \draw[double arrow=black] (iiDEP) -- (MCP);
    \draw[double arrow=black] (MCP) -- (FCP);
    \draw[double arrow=black] (FCP) -- (DCP);
    \draw[double arrow=black] (DCP) -- (LMP);

    \draw[double arrow=black] (CEP) -- (CPP);
    \draw[double arrow=black] (CPP) -- (LMP);
        %\node[vertex, draw=black, fill=white] (HDP) at (2*3.875,13.5+1.5) {\HDP};  
        %\node[vertex, draw=black, fill=white] (JDP) at (2*3.875,10.5) {\JDP};
        %\draw[double arrow=black] (HDP) -- (JDP);
        %\draw[double arrow=black] (JDP) -- (iLTP);
\end{tikzpicture}
\caption{Extension and convergence theorems in measure and integration theory, stated as axioms for an ordered field, are each a sufficient condition for the Lebesgue Integral Property and Lebesgue Measure Property, which were proven in \cite{can24} to be equivalent to Dedekind completeness.}\label{TheFig}
\end{figure}

In accordance with the explanation in the previous section, the diagram in Figure~\ref{TheFig} shows that our strategy is to show that each statement in question is a sufficient condition for \LIP\  and \LMP, which have been proven in \cite[Theorem~1]{can24} to be equivalent to Dedekind completeness. That is, if we succeed in proving every one-directional implication in Figure~\ref{TheFig}, then Dedekind completeness implies an arbitrary statement from Section~\ref{StateSec} which implies one of \LIP, \LMP, which implies Dedekind completeness. Hence, each such statement is an equivalent completeness axiom for $\F$. Since some chains of implications can be formed between statements, we highlighted in gray the strongest statements, which are \iiDEP\  and \CEP, and also the weakest statements, which should be \LIP\  and \LMP, as previously explained. Some of the implications are immediate from the preliminaries as we shall now show.

\section{Implications immediate from the theory}\label{ImpliesSec}

If a positive linear functional $\elemI:\varphi\mapsto\int_a^b\varphi$ on the vector lattice $\elemLatt$ is extendable to a Daniell integral on the second Daniell extension of $\elemLatt$, then the first step of this extension process must be possible, so $\elemI$ has a continuous extension to the first Daniell extension of $\elemLatt$. Necessarily, $\elemI$ in this case must be a Daniell functional, so we have
\[\iiDEP\implies\iDEP\implies\DFP.\]
If the function $E\mapsto\int_a^b\chi_E$ is extendable to a \Cara outer measure on $\abClosed$, then this extension process necessitates that $E\mapsto\int_a^b\chi_E$ is a premeasure, and if it is so, then every element of the algebra of all measurable subsets $E$ of $\abClosed$ must be assigned an integral $\int_a^b\chi_E\in\F$. Thus,
\[\CEP\implies\CPP\implies\LMP.\]
The very essence of the Daniell extension framework for integration theory is that if the extension of a Daniell functional exists for the second Daniell extension of the underlying vector lattice, then the Daniell integral satisfies a Monotone Convergence Principle, which implies an analog of Fatou's lemma, which in turn implies an analog of the Dominated Convergence Principle. See, for instance, \cite[Propositions~10,12,13]{roy88}. The statements of these propositions may be instantiated on the specific integrals and functions in our statements of \iiDEP, \MCP, \FCP\  and \DCP. Thus, the implications
\[\iiDEP\implies\MCP\implies\FCP\implies\DCP,\]
are immediate. As mentioned earlier, the equivalence of \LIP\  and \LMP\  is part of the subject of the paper \cite{can24}. Thus, by inspection of Figure~\ref{TheFig}, only two implications remain, which are ``\DFP\  implies \LIP,'' and ``\DCP\  implies \LMP.''

\section{Summary and final remarks}

To complete the proofs of the implications in Figure~\ref{TheFig}, we state and prove our main result.


\begin{theorem} The statements in Figure~\ref{TheFig} are each equivalent to the Dedekind completeness of the ordered field $\F$.
\end{theorem}

\begin{proof} As discussed in Section~\ref{ImpliesSec}, only two implications remain, which we now prove. Consider an interval $\abClosed\sub\F$.\\

\noindent$\DFP\implies\LIP$. Let $f:\abClosed\into\F$ be a Lebesgue measurable function, and suppose $(\varphi_n)$ and $(\psi_n)$ are a monotonically decreasing sequences of step functions that converge to $f$ almost everywhere on $\abClosed$. The sequence $(\varphi_n\wedge\psi_n)$ also monotonically decreases to $f$, and the sequence $(\varphi_n-\varphi_n\wedge\psi_n)$ monotonically decreases to the zero function. If \DFP\ is true, then we have $\nseqlim\int_a^b(\varphi_n-\varphi_n\wedge\psi_n)=0$, which implies $\nseqlim\int_a^b\varphi_n=\nseqlim\int_a^b(\varphi_n\wedge\psi_n)$. Since $(\varphi_n)$ and $(\psi_n)$ are arbitrary, we may interchange $\varphi_n$ and $\psi_n$ in the previous equation, and we further have
\[\nseqlim\int_a^b\varphi_n=\nseqlim\int_a^b(\varphi_n\wedge\psi_n)=\nseqlim\int_a^b\psi_n.\]
That is, for any two sequences of step functions that monotonically decrease to $f$ almost everywhere on $\abClosed$, the corresponding sequences of integrals converge to the same value. From the definition in \cite[pp.~114--115]{can24}, this means that $f$ has a Lebesgue integral over $\abClosed$, which proves \LIP.\\ 


\noindent$\DCP\implies\LMP$. Let $E$ be a measurable subset of $\abClosed$, and let $(\varphi_n)$ be a sequence of step functions (which are measurable functions) $\abClosed\into\F$ that monotonically decrease to $\chi_E$ almost everywhere on $\abClosed$. Such a sequence is guaranteed to exist by the definition of measurable subset of $\abClosed$ from \cite[pp.~114--115]{can24}. The fact that $(\varphi_n)$ is monotonically decreasing may be used in a routine epsilon argument, valid in the ordered field $\F$, to show that $\chi_E\leq\varphi_n$ for all $n\in\N$. Also for the same reason that $(\varphi_n)$ is monotonically decreasing, we further have $\chi_E\leq \varphi_n\leq \varphi_1$, but since the left-most function in these inequalities is nonnegative, we may replace $\varphi_n$ by $|\varphi_n|$. That is, $|\varphi_n|\leq \varphi_1$ for all $n\in\N$. If \DCP\  holds, then $\int_a^b\chi_E=\nseqlim\int_a^b\varphi_n$. Hence, $\chi_E$ has a Lebesgue integral over $\abClosed$, or equivalently, $E$ has Lebesgue measure, which proves \LMP.
\end{proof}

\begin{thebibliography}{99}
\bibitem[can24]{can24} R. R. S. Cantuba, \textit{Littlewood's principles in reverse real analysis}, Real Anal. Exchange, \textbf{49(1)} (2024), 111--122.
\bibitem[roy88]{roy88} H. L. Royden, \textit{Real Analysis}, Third edition. Macmillan Publishing Company, New York, 1988.
\end{thebibliography}
\end{document}
